\documentclass[10pt, conference, letterpaper]{IEEEtran}
\IEEEoverridecommandlockouts

\usepackage{times}
\usepackage{cite}
\usepackage{amsmath,amssymb,amsfonts}
\usepackage{algorithm}
\usepackage{algorithmic}
\usepackage{graphicx}
\usepackage{textcomp}
\usepackage{booktabs}
\usepackage{multirow}
\usepackage{url}
\usepackage{microtype}
\usepackage{balance}

% Enforce exact IEEE conference layout parameters
\setlength{\columnsep}{0.17in}
\setlength{\parindent}{3.5mm}

\def\BibTeX{{\rm B\kern-.05em{\sc i\kern-.025em b}\kern-.08em
    T\kern-.1667em\lower.7ex\hbox{E}\kern-.125emX}}

\begin{document}

\title{MindScreen: A Decision-Level Multimodal Fusion Architecture
  for Depression Screening Support with Deterministic High-Risk Escalation\\
\thanks{Capstone project, Dept.\ of Computer Science and Engineering,
  RV Institute of Technology and Management, Bengaluru, India.}}

\author{%
\IEEEauthorblockN{Savitha G\IEEEauthorrefmark{1},
  Ruchita Saraf\IEEEauthorrefmark{2},
  Shravya Sanikere\IEEEauthorrefmark{2},
  Chandrika Lamani\IEEEauthorrefmark{2},
  Apoorva K\IEEEauthorrefmark{2}}
\IEEEauthorblockA{\textit{Dept.\ of Computer Science and Engineering,
  RV Institute of Technology and Management, Bengaluru, India}\\
  \footnotesize
  \IEEEauthorrefmark{1}savithag.rvitm@rvei.edu.in \quad
  \IEEEauthorrefmark{2}\{ruchitasaraf9, shravya.sanikere,
    chandrika.lamani, apoorva.k\}@gmail.com}}

\maketitle
\thispagestyle{empty}
\pagestyle{empty}

% ─────────────────────────────────────────────────────────────────────────────
\begin{abstract}
This paper presents \textbf{MindScreen}, a full-stack mental-health screening
support engineering prototype that fuses three
complementary signal sources---Patient Health Questionnaire-9 (PHQ-9) responses,
transformer-based text emotion classification, and client-side Web Audio
acoustic feature extraction (mean RMS loudness, energy variability, zero-crossing
rate, spectral centroid, spectral rolloff, and speaking ratio)---through
decision-level late fusion.
A key architectural contribution is the \textbf{High-Risk Escalation (HRE)}
module: a formally specified deterministic layer that prevents high-risk
indicators (PHQ-9 Item~9 endorsement, aggregate score $S \ge 20$, or affirmative
crisis intent detected by a scope-aware rule-based linguistic filter) from being
diluted by modality averaging.
The text branch integrates deterministic first-person, negation, attribution,
quotation, and idiom rules as an initial crisis-language safety gate.
The prototype integrates \textbf{Saathi}, an LLM-assisted conversational
companion backed by Google Gemini~2.5 Flash with crisis-resource routing.
Evaluation comprises: (i)~HRE anti-dilution verification and 5-way weight
sensitivity analysis across five constructed evaluation profiles;
(ii)~scope-rule verification on a constructed 17-statement corpus, on which the
implemented detector matched all 17 team-assigned labels; and (iii)~in-process
latency profiling of deterministic local functions. Hosted-API and end-to-end
deployment latency were not measured.
The system is explicitly framed as an engineering screening support prototype.
No clinical validation or diagnostic accuracy has been established; both are
identified as primary future work.
\end{abstract}

\begin{IEEEkeywords}
Multimodal late fusion, depression screening, PHQ-9, transformer NLP,
rule-based escalation, affective computing, mental health support.
\end{IEEEkeywords}

% ─────────────────────────────────────────────────────────────────────────────
\section{Introduction}
\label{sec:intro}

Major Depressive Disorder (MDD) affects over 280 million people globally and
is the leading cause of non-fatal disability~\cite{who2023}.
In India, the mental healthcare treatment gap exceeds 75\%--85\% due to a
shortage of practitioners ($\approx$0.75 per 100{,}000 population) and
entrenched social stigma~\cite{gururaj2016}.
While population screening instruments such as the K6~\cite{kessler2003} and
PHQ-9~\cite{kroenke2001} are widely used and validated, self-reports alone
cannot capture real-time affective language or vocal cues that may indicate distress.

Automated multimodal screening tools address this gap, but unconstrained
probabilistic averaging introduces a critical failure mode:
\textit{false-negative dilution}.
If a user reports acute suicidal ideation on PHQ-9 Item~9 or writes crisis
language in free text, but exhibits calm vocal acoustics, linear weighted
averaging can push the composite risk below the high-risk threshold,
potentially reducing the fused score below the selected high-priority
threshold and thereby failing to preserve the explicit escalation signal.
Among the systems and descriptions reviewed for this study, we did not identify
an explicitly documented deterministic high-risk override mechanism matching
the HRE design used in MindScreen. This observation is limited by the scope of
the reviewed literature and publicly available system descriptions.

The engineering contributions of this work are:
\begin{enumerate}
  \item A \textbf{tri-modal late-fusion pipeline} over PHQ-9, text, and
    audio signal sources;
  \item A \textbf{High-Risk Escalation (HRE) module}---a deterministic
    rule layer that bypasses fusion when explicit high-risk indicators are
    present, with formal mathematical specification and verification;
  \item A \textbf{conversational companion (Saathi)} with LLM-backed
    empathetic dialogue and crisis-resource routing to national helplines; and
  \item A \textbf{reproducible local evidence script} covering HRE semantics,
    the constructed crisis-language corpus, and in-process deterministic-function
    latency with raw observations and Python-managed heap measurements.
\end{enumerate}

The system is presented as a screening support prototype.
Its outputs are risk indicators to inform timely clinician consultation, not
diagnostic labels.
No clinical validation has been conducted; this is explicitly acknowledged
throughout.

% ─────────────────────────────────────────────────────────────────────────────
\section{Related Work}
\label{sec:related}

Existing mental-health screening systems fall into three paradigms.
\textit{Text-only systems}~\cite{ji2022,xu2024} achieve strong benchmark
metrics but ignore somatic signals.
\textit{Acoustic systems}: several acoustic depression-assessment
studies~\cite{cummins2015,scherer2013} use recordings collected under controlled
or semi-controlled conditions, which may limit direct transfer to unconstrained
consumer-device environments.
\textit{Multimodal systems}~\cite{yang2023,ringeval2019,devault2014} predominantly
use early-feature concatenation, which is sensitive to modality missingness and
produces opaque intermediate representations.
Conversational screening tools: the cited Woebot study~\cite{fitzpatrick2017}
evaluated a text-based conversational agent in a randomized controlled trial;
its task, population, and evaluation setting differ from the multimodal
screening architecture investigated here.

Table~\ref{tab:related} positions MindScreen against these approaches.
Two distinguishing design properties are: (a)~decision-level late fusion that
preserves per-modality interpretability, and (b)~a formally specified HRE
module designed to provide deterministic high-risk override behavior.
Within the systems and descriptions reviewed for this study, we did not identify
an explicitly documented deterministic high-risk override mechanism matching the
HRE design used in MindScreen.

\begin{table*}[!t]
\caption{Comparison of Related Mental Health Screening Systems}
\label{tab:related}
\centering
\footnotesize
\setlength{\tabcolsep}{8pt}
\begin{tabular}{@{}lcccc@{}}
\toprule
\textbf{System} & \textbf{Modalities}
  & \textbf{Fusion} & \textbf{HRE} & \textbf{Technical Hosting} \\
\midrule
Woebot~\cite{fitzpatrick2017}        & Text       & —          & Static & Cloud service \\
MentalBERT~\cite{ji2022}             & Text       & —          & No     & Pretrained model \\
SimSensei~\cite{devault2014}         & Text+V+A   & Early      & No     & Kiosk system \\
AVEC~\cite{ringeval2019}             & Text+V+A   & Early/mid  & No     & Benchmark dataset \\
Cummins~et~al.~\cite{cummins2015}   & Audio      & —          & No     & Research pipeline \\
\textbf{MindScreen (this work)}      & \textbf{PHQ+Text+A} & \textbf{Late} & \textbf{Yes} & \textbf{Engineering prototype} \\
\bottomrule
\end{tabular}
\vspace{1ex}

\parbox{0.90\textwidth}{\footnotesize \textit{Note:} The systems listed in
Table~\ref{tab:related} are not directly comparable in terms of task definition,
target population, clinical dataset, deployment status, or evaluation protocol;
the table is intended to provide architectural context rather than a controlled
performance comparison. Repository deployment configuration demonstrates a
hostable prototype, but no controlled cloud performance environment was verified.}
\end{table*}

% ─────────────────────────────────────────────────────────────────────────────
\section{System Architecture}
\label{sec:method}

\subsection{Input Space and Output Taxonomy}

An assessment session accepts a tuple
$\mathcal{X} = (\mathbf{q},\, T,\, \mathbf{a})$ where
$\mathbf{q} = (q_1, \dots, q_9)$ with $q_i \in \{0,1,2,3\}$ are the PHQ-9 item
responses, $T \in \mathcal{V}^*$ is a free-text journal entry, and
$\mathbf{a}$ is the captured audio signal.
The system outputs an internal screening category $\hat{y}$ drawn from the
4-tier set:
\begin{equation}
  \mathcal{C} = \{c_0\!:\!\text{minimal},\;c_1\!:\!\text{mild},\;
                  c_2\!:\!\text{moderate},\;c_3\!:\!\text{high-priority}\},
  \label{eq:label_space}
\end{equation}
together with an output confidence/priority score $c_{\text{out}} \in [0,1]$, a
Boolean crisis indicator $C_F$, and a resource display flag $R_D$.

The $c_3$ category is an internal high-priority screening tier and should not be
interpreted as a direct diagnosis of severe depression or acute suicide risk.
For backward compatibility, the current API serializes this tier using the
legacy string ``severe''; the user interface presents it as ``High Priority.''
MindScreen maps PHQ-9 scores to this 4-tier space by grouping the clinically
distinct moderately-severe (15--19) and severe (20--27) PHQ-9 score bands~\cite{kroenke2001}
into this single high-priority tier ($c_3$, Table~\ref{tab:phq_map}).
Both bands trigger professional-referral guidance and crisis-resource display
in the implemented system, motivating this operational grouping.
This internal category mapping is an engineering screening threshold and
must not be conflated with the formal clinical cutoff for severe depression
($S \ge 20$).

\begin{table}[!t]
\caption{PHQ-9 Score to MindScreen Internal Screening Tier Mapping}
\label{tab:phq_map}
\centering
\resizebox{\columnwidth}{!}{%
\begin{tabular}{@{}clll@{}}
\toprule
\textbf{Score $S$} & \textbf{PHQ-9 Band~\cite{kroenke2001}} & \textbf{Internal Tier ($\hat{y}$)}
  & \textbf{Operational Action} \\
\midrule
0--4   & Minimal           & $c_0$ (Minimal)       & Routine self-monitoring \\
5--9   & Mild              & $c_1$ (Mild)          & Wellness exercises \\
10--14 & Moderate          & $c_2$ (Moderate)      & Recommend consultation \\
15--19 & Moderately severe & \multirow{2}{*}{$c_3$ (High-priority)}
  & \multirow{2}{*}{\parbox{2.6cm}{\raggedright Professional referral;\\display helplines ($R_D$)}} \\
20--27 & Severe            &          & \\
\bottomrule
\end{tabular}%
}
\end{table}

\subsection{Modality Processing Pipelines}

\subsubsection{PHQ-9 branch ($\mathbf{p}^Q$)}
The aggregate score $S = \sum_{i=1}^{9} q_i$ is mapped to a fixed 4-class
score vector via a piecewise function $\Phi_Q$:
\begin{equation}
  \mathbf{p}^Q =
  \begin{cases}
    [0.80,\;0.15,\;0.05,\;0.00]^T, & 0  \le S \le 4  \\
    [0.10,\;0.70,\;0.15,\;0.05]^T, & 5  \le S \le 9  \\
    [0.00,\;0.15,\;0.70,\;0.15]^T, & 10 \le S \le 14 \\
    [0.00,\;0.05,\;0.15,\;0.80]^T, & 15 \le S \le 27
  \end{cases}
  \label{eq:phq_probs}
\end{equation}
These probability vectors are manually specified engineering constants aligned
with conventional PHQ-9 score intervals~\cite{kroenke2001}; they represent heuristic design
parameters rather than empirical class-conditional likelihoods estimated from
patient data.
The PHQ-9 branch uses band-level heuristic vectors rather than a
score-continuous probabilistic model. Consequently, two respondents with
different scores within the same interval receive the same questionnaire
branch vector. This design choice prioritizes deterministic implementation
and interpretability over fine-grained score calibration.

\subsubsection{NLP text branch ($\mathbf{p}^T$)}
Free-text input $T$ is classified by \textit{j-hartmann/emotion-english-distilroberta-base}~\cite{hartmann2022},
accessed via the HuggingFace Inference API.
The top emotion class $e^*$ and its maximum model output score $s$ are extracted:
\begin{equation}
  e^* = \arg\max_{e \in \mathcal{E}} P(e \mid T;\,\Theta_{\text{NLP}}),
  \qquad
  s = P(e^* \mid T;\,\Theta_{\text{NLP}}).
  \label{eq:nlp_argmax}
\end{equation}
A heuristic mapping $\Gamma: \mathcal{E} \to \{0,1,2,3\}$ translates the
7-class emotion output to a screening index:
\begin{equation}
  \Gamma(e^*) =
  \begin{cases}
    0, & e^* \in \{\text{joy},\;\text{neutral},\;\text{surprise}\} \\
    1, & e^* \in \{\text{fear},\;\text{anger}\} \\
    2, & e^* \in \{\text{sadness},\;\text{disgust}\}
  \end{cases}
  \label{eq:emotion_map}
\end{equation}
If affirmative crisis intent is detected in $T$ by the deterministic safety gate,
the index is escalated: $\Gamma(e^*) \leftarrow 3$. The gate matches a bounded
set of first-person or explicitly self-directed multi-word expressions. Predicate-
attached negation rules suppress statements such as ``I do not want to die'',
while third-party speech, quotations, and idioms are handled separately. Idiom
spans are masked rather than causing early termination so a later first-person
clause remains detectable. General distress words such as ``hopeless'' and
``worthless'' do not by themselves assert acute crisis intent. This deterministic
gate is not a clinically validated suicide-risk detector and cannot cover all
indirect, multilingual, metaphorical, or culturally specific expressions.
The score vector $\mathbf{p}^T$ is then formed by assigning $s$ to the
predicted class and $0.10$ to each remaining class, followed by
$\ell_1$ normalization.
The text branch is used as a coarse affective signal rather than a standalone
depression classifier. The mapping from emotion categories to screening indices
is a heuristic design choice; it is not intended to establish diagnostic severity
or to infer the absence of crisis risk from non-negative emotion labels.
When the HuggingFace API is unreachable or times out, a deterministic keyword-rule
fallback maintains software-level service continuity; this heuristic fallback
provides operational availability rather than equivalent transformer classification fidelity.

\subsubsection{Audio branch ($\mathbf{p}^A$)}
To avoid adding server-side acoustic processing dependencies, MindScreen
extracts signal descriptors in the client browser via the standard Web Audio
API. No deployment-memory comparison was measured for this design choice.
During voice input, the browser captures audio via an \texttt{AudioContext}
(native browser sampling at 44.1 or 48~kHz).
An \texttt{AnalyserNode} with FFT size $N_{\text{fft}} = 2048$ computes real-time
frame spectra. Frames with RMS below a silence threshold of $0.02$ are classified
as unvoiced silence.
Across all active frames, six acoustic descriptors are accumulated:
mean RMS loudness ($\text{RMS}_\mu$), RMS energy variability ($\text{RMS}_\sigma$),
zero-crossing rate ($\text{ZCR}$), normalized spectral centroid ($S_C$),
85\% spectral rolloff ($S_R$), and speaking-to-silence ratio ($R_{\text{speak}}$).
The acoustic coefficients and threshold intervals were selected as engineering
design parameters for prototype-level signal aggregation, rather than estimated
through supervised learning. Their purpose is to provide a lightweight and
interpretable heuristic representation of acoustic variation. The selected
features and coefficient values should therefore not be interpreted as validated
biomarkers of depression. Future work will estimate the mapping using an
appropriately annotated speech dataset and evaluate its robustness across
recording devices, languages, accents, and speaking conditions.

The normalized feature vector
$\mathbf{f} = [\text{RMS}_\mu, \text{RMS}_\sigma, \text{ZCR}, S_C, S_R, R_{\text{speak}}]^T \in [0,1]^6$
is transmitted to the server, where a linear composite depression index
$I_D \in [0,1]$ is computed:
\begin{equation}
  I_D = \sum_{j=1}^6 \alpha_j \, \phi_j(\mathbf{f}),
  \label{eq:acoustic_index}
\end{equation}
where $\phi_j$ measures deviation toward depressive characteristics:
$\phi_1 = 1 - \text{RMS}_\mu$ (reduced vocal energy),
$\phi_2 = 1 - \min(\text{RMS}_\sigma/0.25, 1)$ (reduced variability / flat affect),
$\phi_3 = 1 - S_C$ (low spectral centroid / muffled vocal tone),
$\phi_4 = 1 - S_R$ (suppressed high-frequency spectral rolloff),
$\phi_5 = 1 - R_{\text{speak}}$ (reduced speaking ratio / hesitancy), and
$\phi_6 = \min(|\text{ZCR} - 0.25|/0.25, 1)$ (deviation from typical voiced baseline),
weighted by design constants
$\boldsymbol{\alpha} = [0.20, 0.20, 0.15, 0.15, 0.20, 0.10]^T$ ($\sum \alpha_j = 1$).
The index $I_D$ is mapped to the 4-class score vector $\mathbf{p}^A$ via
piecewise linear interpolation between boundary anchors:
\begin{equation}
  \mathbf{p}^A(I_D) = (1-t)\,\mathbf{a}_k + t\,\mathbf{a}_{k+1}, \quad t = \frac{I_D - l_k}{u_k - l_k},
  \label{eq:audio_interp}
\end{equation}
evaluated across four intervals $[l_k, u_k)$: $[0.00, 0.25)$ ($k=0$),
$[0.25, 0.45)$ ($k=1$), $[0.45, 0.65)$ ($k=2$), and $[0.65, 1.00]$ ($k=3$),
with design anchors $\mathbf{a}_0 = [0.80, 0.15, 0.04, 0.01]^T$,
$\mathbf{a}_1 = [0.50, 0.35, 0.12, 0.03]^T$,
$\mathbf{a}_2 = [0.15, 0.55, 0.25, 0.05]^T$,
$\mathbf{a}_3 = [0.05, 0.20, 0.60, 0.15]^T$, and
$\mathbf{a}_4 = [0.02, 0.08, 0.35, 0.55]^T$, followed by $\ell_1$ normalization.
If voice input is omitted by the user, a conservative prior vector
$\mathbf{p}^A = [0.25, 0.45, 0.20, 0.10]^T$ is assigned.
This distribution skews probability mass slightly toward the mild tier ($c_1$)
as an intentional engineering safeguard, preventing unobserved audio from
prematurely asserting an absence of distress. The client sends a null audio
field, the server reports audio as unavailable, and the fixed 30\% audio weight
is applied to this prior. No acoustic observation is fabricated, weights are
not renormalized, and this missing-modality policy has not been statistically validated.

\subsection{Decision-Level Late Fusion and Score Transformation}
\label{sec:fusion}

The three score vectors are combined as a convex sum:
\begin{equation}
  \mathbf{p}^F = w_T\,\mathbf{p}^T + w_A\,\mathbf{p}^A + w_Q\,\mathbf{p}^Q,
  \label{eq:fusion}
\end{equation}
with weights $w_T = 0.50$, $w_A = 0.30$, $w_Q = 0.20$ ($\sum w_m = 1$).
The selection of these weights is an engineering design choice grounded in the
complementary characteristics of the three channels:
(1)~the text channel ($w_T = 0.50$) carries direct linguistic expressions of affect,
cognitive distortion, and acute intent, warranting primary weight;
(2)~the audio channel ($w_A = 0.30$) captures real-time somatic, prosodic, and psychomotor markers
(such as flat affect and hesitant speech); and
(3)~the PHQ-9 channel ($w_Q = 0.20$) provides a retrospective two-week psychometric baseline.
Because each $\mathbf{p}^m$ sums to one and the weights sum to one,
$\mathbf{p}^F$ already lies on the simplex; normalization is retained
defensively.
Following convex combination, the distribution is passed through a
temperature-scaled score transformation~\cite{guo2017} ($T = 1.20$) to soften
heuristic overconfidence:
\begin{equation}
  p_k^{\text{soft}} = \frac{\exp(\log(\max(p_k^F, \epsilon))/T)}{\sum_{j=1}^4 \exp(\log(\max(p_j^F, \epsilon))/T)},
  \label{eq:temperature}
\end{equation}
where $\epsilon = 10^{-7}$ provides numerical stability against underflow.
The temperature parameter ($T = 1.20$) was selected to counteract heuristic probability saturation.
Because individual branch vectors contain sharp heuristic peaks (e.g., $0.80$--$0.90$), direct convex
combination can produce overly confident fused scores ($>0.85$) prior to safety verification.
Setting $T = 1.20$ provides moderate entropy expansion, softening peaked probabilities into a
smoother distribution without altering their ordinal ranking.
We emphasize that this parameter ($T = 1.20$) represents an experimental
heuristic score-softening transformation to attenuate extreme probabilities;
it is not a statistically calibrated posterior probability obtained via
validation-set optimization (such as Platt scaling or isotonic regression).
The temperature-softened score vector $\mathbf{p}^{\text{soft}}$ is used to obtain the base
screening category and base score before escalation:
\begin{equation}
  \hat{y} = \arg\max_k p_k^{\text{soft}},
  \qquad
  s_F = \max_k p_k^{\text{soft}}.
  \label{eq:prediction}
\end{equation}

\subsection{High-Risk Escalation (HRE) Module}
\label{sec:hre}

Decision-level fusion still exposes the false-negative dilution problem: a
user who scores $S = 22$ on the PHQ-9 but generates joy-emotion text and
active audio will produce a fused vector dominated by the text and audio
channels, potentially outputting ``mild.''
The HRE module resolves this by applying two deterministic rules
\textit{after} fusion, unconditionally overriding the fused output when
explicit high-risk indicators are present.

To ensure semantic clarity, the system maintains a strict separation among:
(1)~the internal screening category $\hat{y} \in \{c_0, c_1, c_2, c_3\}$;
(2)~the Boolean crisis indicator $C_F$, which is asserted specifically upon
detecting acute suicidal ideation ($q_9 > 0$ or affirmative crisis language);
(3)~the resource display flag $R_D$, which presents crisis helpline information
(Tele-MANAS~\cite{telemanas2022}, KIRAN~\cite{kiran2020}) whenever $C_F$ is active or as a
precautionary measure whenever $\hat{y} = c_3$; and
(4)~the priority score $c_{\text{out}}$, which is clamped to a floor of $0.90$
under rule-based override.
Resource display is intentionally precautionary: the presentation of crisis-related
information does not imply that the system has established suicidal intent or
a clinical diagnosis. The resource-routing behavior is an operational safety
design choice.
The value $c_{\text{out}} \leftarrow \max(0.90, s_F)$ is an engineering priority
indicator intended for triage sorting, not a calibrated posterior probability.

\noindent\textbf{Rule~1 — High-score override.} If $S \ge 20$:
\begin{equation}
  \hat{y} \leftarrow c_3,\quad
  c_{\text{out}} \leftarrow \max(0.90,\; s_F),\quad
  R_D \leftarrow \text{TRUE}.
  \label{eq:hre_severe}
\end{equation}

\noindent\textbf{Rule~2 — Explicit-indicator escalation.}
If $q_9 > 0$ (suicidal ideation endorsed on PHQ-9 Item~9) or affirmative crisis
intent is detected in $T$:
\begin{equation}
  \hat{y} \leftarrow c_3,\quad
  c_{\text{out}} \leftarrow \max(0.90,\; s_F),\quad
  C_F \leftarrow \text{TRUE},\quad
  R_D \leftarrow \text{TRUE}.
  \label{eq:hre_escalate}
\end{equation}
Table~\ref{tab:hre} enumerates the complete HRE decision matrix.

\begin{table}[!t]
\caption{HRE Module: Complete Decision Matrix}
\label{tab:hre}
\centering
\resizebox{\columnwidth}{!}{%
\begin{tabular}{@{}lcccc@{}}
\toprule
\textbf{Trigger Condition} & $\hat{y}$ & $C_F$ & $R_D$ & $c_{\text{out}}$ \\
\midrule
No trigger (fused output)  & \eqref{eq:prediction} & F & F & $s_F$ \\
Aggregate score $S \ge 20$  & $c_3$ & F$^\dagger$ & T & $\max(0.90, s_F)$ \\
PHQ-9 Item 9 ($q_9 > 0$)   & $c_3$ & T & T & $\max(0.90, s_F)$ \\
Crisis language in $T$     & $c_3$ & T & T & $\max(0.90, s_F)$ \\
Fused $\hat{y} = c_3$ alone& $c_3$ & F & T & $s_F$ \\
Multiple triggers          & $c_3$ & T & T & $\max(0.90, s_F)$ \\
\bottomrule
\multicolumn{5}{@{}l}{\footnotesize $^\dagger C_F$ remains FALSE unless Item 9 or crisis language is also present.}
\end{tabular}%
}
\end{table}

Algorithm~\ref{alg:fusion} gives the complete inference flow.

\begin{algorithm}[!t]
\caption{MindScreen Tri-Modal Inference with HRE}
\label{alg:fusion}
\begin{algorithmic}[1]
\REQUIRE $\mathbf{q} \in \{0,1,2,3\}^9$, text $T$, optional audio features $\mathbf{f} \in [0,1]^6$
\ENSURE Tier $\hat{y}$, score $c_{\text{out}}$, crisis flag $C_F$, resource flag $R_D$
\STATE $\mathbf{p}^Q \leftarrow \Phi_Q\!\left(\sum_{i=1}^9 q_i\right)$
\STATE $\mathbf{p}^T \leftarrow \text{NLP}(T)$;\; $\mathbf{p}^A \leftarrow \text{Audio}(\mathbf{f})$ if supplied, otherwise the fixed prior
\STATE $\mathbf{p}^F \leftarrow 0.50\,\mathbf{p}^T + 0.30\,\mathbf{p}^A + 0.20\,\mathbf{p}^Q$
\STATE Compute softened vector $\mathbf{p}^{\text{soft}}$ via \eqref{eq:temperature} ($T=1.20$)
\STATE $\hat{y} \leftarrow \arg\max_k p_k^{\text{soft}}$;\; $s_F \leftarrow \max_k p_k^{\text{soft}}$;\; $c_{\text{out}} \leftarrow s_F$
\STATE $C_F \leftarrow \text{FALSE}$;\; $R_D \leftarrow (\hat{y} == c_3)$
\IF{$S \ge 20$}
  \STATE $\hat{y} \leftarrow c_3$;\; $c_{\text{out}} \leftarrow \max(0.90, c_{\text{out}})$;\; $R_D \leftarrow \text{TRUE}$
\ENDIF
\IF{$q_9 > 0 \lor \text{CrisisIntent}(T)$}
  \STATE $\hat{y} \leftarrow c_3$;\; $c_{\text{out}} \leftarrow \max(0.90, c_{\text{out}})$
  \STATE $C_F \leftarrow \text{TRUE}$;\; $R_D \leftarrow \text{TRUE}$
\ENDIF
\RETURN $\hat{y},\; c_{\text{out}},\; C_F,\; R_D$
\end{algorithmic}
\end{algorithm}

\subsection{Saathi: LLM-Assisted Conversational Companion}
\label{sec:saathi}

Saathi is an interactive, LLM-assisted conversational component backed by
Google Gemini~2.5 Flash, maintaining an 8-turn conversational context.
Saathi is invoked post-screening as a distinct conversational companion and is
not part of the deterministic local screening benchmark.
It operates in four sequential stages:
(1)~a \textit{crisis safety gate} utilizing linguistic negation and
scope-resolved intent analysis to immediately surface national helpline
information if affirmative crisis signals appear; the Saathi crisis gate is a
rule-based first-stage safeguard and does not guarantee detection of all
expressions of suicidal intent or other forms of acute psychological risk;
(2)~\textit{affective topic routing} classifying user statements into six
domains (family, academics, insomnia, anxiety, isolation, positive);
(3)~\textit{LLM response generation} conditioned on the classified domain and
conversation history; and
(4)~\textit{structured intervention suggestions} (4-7-8 breathing, 5-4-3-2-1
grounding) interleaved every 3--4 turns.
Saathi does not provide clinical advice; it directs users to professional
resources.
LLM response safety and factuality have not been formally evaluated.
Local template fallbacks maintain service continuity if external API quotas
are exhausted.

\subsection{Privacy, Consent, and Data Governance}
\label{sec:privacy}

Given the sensitive nature of psychological screening inputs, MindScreen implements
four core data-governance principles:
\begin{enumerate}
  \item \textbf{Informed user consent}: Prior to initiating screening, users review
    an explicit consent disclaimer stating that MindScreen is an engineering research
    prototype, not a diagnostic or emergency medical device, and does not replace professional
    consultation.
  \item \textbf{Client-side audio processing}: The browser creates an in-memory
    recording for playback and extracts six numerical descriptors. The application
    does not upload the raw audio blob; only the descriptor vector is sent.
  \item \textbf{External API limitation}: Free-text journal and companion inputs
    can be sent to Hugging Face or Google Gemini endpoints. The prototype does not
    implement PII detection or redaction, so users must not assume entered PII is removed.
  \item \textbf{Local inference fallback}: If hosted inference fails, deterministic
    local rules provide a reduced-fidelity screening path after the attempted request.
\end{enumerate}

% ─────────────────────────────────────────────────────────────────────────────
\section{Evaluation}
\label{sec:eval}

\subsection{HRE Rule Execution Verification and Weight Sensitivity}

We executed four isolated scenarios through the production fusion, score
transformation, HRE, audio, PHQ-9, and crisis-gate functions using a fixed local
text score vector to avoid hosted-API variability (Table~\ref{tab:hre_verify}).
HRE-1 isolates $S\ge20$ with $q_9=0$ and no crisis language; it escalates the
tier and resources while leaving $C_F$ false. HRE-2 isolates Item~9, HRE-3
isolates affirmative text, and HRE-4 verifies the no-trigger path.

\begin{table}[!t]
\caption{HRE Verification: Adversarial Masking Scenarios}
\label{tab:hre_verify}
\centering
\resizebox{\columnwidth}{!}{%
\begin{tabular}{@{}p{2.2cm}p{2.4cm}ccc@{}}
\toprule
\textbf{Scenario} & \textbf{Signal State}
  & \textbf{Tier} & $C_F$ & $R_D$ \\
\midrule
HRE-1: High score only & $S=20,q_9=0$, ordinary text
  & $c_3$ ($0.90$) & F & T \\
HRE-2: Item~9 & $S=1,q_9=1$, ordinary text
  & $c_3$ ($0.90$) & T & T \\
HRE-3: Crisis text & $S=0$, ``I want to die''
  & $c_3$ ($0.90$) & T & T \\
HRE-4: No trigger & $S=0,q_9=0$, ordinary text
  & $c_0$ ($0.4907$) & F & F \\
\bottomrule
\end{tabular}%
}
\end{table}

The score distributions remain base heuristic distributions; the 0.90 value is
an escalation priority floor rather than a probability. These constructed cases
verify software semantics only and do not establish clinical performance.

To evaluate the sensitivity of late fusion to weighting choices and test the
anti-dilution behavior of the HRE module, we conducted an empirical sensitivity
analysis across five weighting configurations: Equal (33/33/33), Text-Dominant
(60/20/20), Audio-Dominant (20/60/20), Questionnaire-Dominant (20/20/60), and
MindScreen (50/30/20).
Under concordant high symptom burden and minimal profiles, all five schemes converged
consistently (high symptom burden: $0.56$--$0.75$, minimal: $0.65$--$0.80$).
Under masked affect ($S=20$, positive text), all linear schemes assigned
minimal or mild categories ($0.37$--$0.60$ score), and were rescued only by the
deterministic HRE override.
Under implicit ideation ($S=2, q_9=1$), questionnaire-dominant weighting
assigned minimal risk tier ($c_0$ with score $0.53$) despite Item~9 endorsement,
which the HRE module rescued to $c_3$ ($c_{\text{out}} = 0.90$, Table~\ref{tab:ablation}).
These constructed scenarios illustrate that the evaluated linear weighting
configurations can be vulnerable to discordant inputs, motivating the use of a
deterministic HRE layer.

For reproducibility, each constructed profile was evaluated using explicitly
specified modality score vectors ($\mathbf{p}^Q, \mathbf{p}^T, \mathbf{p}^A$):
•~P1 ($S=22$): $\mathbf{p}^Q=[0.00, 0.05, 0.15, 0.80]^T$, $\mathbf{p}^T=[0.05, 0.05, 0.05, 0.85]^T$, $\mathbf{p}^A=[0.03, 0.13, 0.46, 0.38]^T$;
•~P2 ($S=1$): $\mathbf{p}^Q=[0.80, 0.15, 0.05, 0.00]^T$, $\mathbf{p}^T=[0.90, 0.05, 0.03, 0.02]^T$, $\mathbf{p}^A=[0.52, 0.34, 0.11, 0.03]^T$;
•~P3 ($S=20$): $\mathbf{p}^Q=[0.00, 0.05, 0.15, 0.80]^T$, $\mathbf{p}^T=[0.90, 0.05, 0.03, 0.02]^T$, $\mathbf{p}^A=[0.32, 0.45, 0.19, 0.04]^T$;
•~P4 ($S=6$): $\mathbf{p}^Q=[0.10, 0.70, 0.15, 0.05]^T$, $\mathbf{p}^T=[0.90, 0.05, 0.03, 0.02]^T$, $\mathbf{p}^A=[0.03, 0.12, 0.43, 0.42]^T$;
•~P5 ($S=2, q_9=1$): $\mathbf{p}^Q=[0.80, 0.15, 0.05, 0.00]^T$, $\mathbf{p}^T=[0.10, 0.15, 0.70, 0.05]^T$, $\mathbf{p}^A=[0.13, 0.48, 0.32, 0.07]^T$.
The scores reported in Table~\ref{tab:ablation} represent $\max_k p_k^F$ before
the temperature transformation and deterministic HRE escalation. The corresponding HRE output
is shown separately using the notation $c_3^*$. All values are rounded to two decimal places.

\begin{table*}[!t]
\caption{Modality Weight Sensitivity and Anti-Dilution Analysis}
\label{tab:ablation}
\centering
\footnotesize
\setlength{\tabcolsep}{6pt}
\begin{tabular}{@{}lccccc@{}}
\toprule
\textbf{Constructed Profile} & \textbf{Equal (33/33/33)} & \textbf{Text-Dom (60/20/20)} & \textbf{Audio-Dom (20/60/20)} & \textbf{PHQ-Dom (20/20/60)} & \textbf{MindScreen (50/30/20)} \\
\midrule
P1: High Symptom Burden ($S \ge 20$) & $0.68 \to c_3^*$ & $0.75 \to c_3^*$ & $0.56 \to c_3^*$ & $0.73 \to c_3^*$ & $0.70 \to c_3^*$ \\
P2: Concordant Minimal               & $0.74$           & $0.80$           & $0.65$           & $0.76$           & $0.77$           \\
P3: Masked Affect                    & $0.41 \to c_3^*$ & $0.60 \to c_3^*$ & $0.37 \to c_3^*$ & $0.49 \to c_3^*$ & $0.55 \to c_3^*$ \\
P4: Somatic Flatness                 & $0.34$           & $0.57$           & $0.30$           & $0.45$           & $0.48$           \\
P5: Implicit Ideation                & $0.36 \to c_3^*$ & $0.49 \to c_3^*$ & $0.35 \to c_3^*$ & $0.53 \to c_3^*$ & $0.46 \to c_3^*$ \\
\bottomrule
\multicolumn{6}{@{}l}{\footnotesize $^*$Escalated to internal high-priority tier ($c_3$) by HRE deterministic bound ($c_{\text{out}} \leftarrow \max(0.90, s_F)$).}
\end{tabular}
\end{table*}

\subsection{Deterministic Crisis-Language Rule Verification}

The safety gate was evaluated on a 17-statement constructed rule-verification
corpus authored by the project team. It includes eight affirmative cases and
nine negative cases spanning direct first-person intent, predicate-attached
negation, third-party speech, quotation, idioms, mixed clauses, and unrelated
negation. The corpus and per-statement output are archived in
\texttt{benchmarks/submission\_evidence.json}. It is not a clinical dataset.

\begin{table}[!t]
\caption{Constructed Crisis-Language Rule Verification}
\label{tab:negation}
\centering
\resizebox{\columnwidth}{!}{%
\begin{tabular}{@{}lcc@{}}
\toprule
\textbf{Metric} & \textbf{Naive substring} & \textbf{Deterministic gate} \\
\midrule
Affirmative detection & 5/8 (62.5\%) & \textbf{8/8 (100.0\%)} \\
False-positive rate & 9/9 (100.0\%) & \textbf{0/9 (0.0\%)} \\
Overall agreement & 5/17 (29.4\%) & \textbf{17/17 (100.0\%)} \\
\bottomrule
\multicolumn{3}{@{}l}{\footnotesize Team-assigned labels on a constructed software-rule corpus only.}
\end{tabular}%
}
\end{table}

The exact-set result verifies the implemented rules for these examples only.
It must not be interpreted as clinical sensitivity, specificity, diagnostic
accuracy, or evidence of generalization to unseen language.
\subsection{Local Deterministic-Function Latency}
\label{sec:latency}

The reproducible evidence script timed 50 invocations after a 10-call warm-up
using \texttt{time.perf\_counter} on Windows~11, AMD64, Python~3.12.3
(Intel64 Family~6 Model~183). Peak values use \texttt{tracemalloc} and therefore
represent Python-managed allocations rather than process RSS. Table~\ref{tab:latency}
covers in-process deterministic functions only; it excludes HTTP handling,
database access, browser feature extraction, hosted APIs, and network round trips.

\begin{table}[!t]
\caption{Local Deterministic-Function Microbenchmark}
\label{tab:latency}
\centering
\resizebox{\columnwidth}{!}{%
\begin{tabular}{@{}lccc@{}}
\toprule
\textbf{Function} & \textbf{Mean (ms)} & \textbf{P95 (ms)} & \textbf{Peak heap (KB)} \\
\midrule
PHQ-9 mapping & 0.0136 & 0.0146 & 1.848 \\
Acoustic scoring & 0.0317 & 0.0407 & 4.445 \\
Crisis-language gate & 0.0496 & 0.0580 & 2.254 \\
Local fusion pipeline & 0.1889 & 0.2083 & 10.021 \\
\bottomrule
\end{tabular}%
}
\end{table}

Raw per-invocation observations and environment metadata are stored in
\texttt{benchmarks/submission\_evidence.json}. These measurements are not
deployment latency. Hosted Hugging Face, Gemini, Render-container, concurrent-load,
and end-to-end cloud latency were not measured and no numerical claims are made
for them.
\section{Limitations and Future Work}
\label{sec:limitations}

\textbf{No clinical validation.}
MindScreen has not been evaluated against a labeled clinical dataset.
No diagnostic accuracy, F-score, AUROC, or formal clinical calibration metrics
are reported.
The system is an adjunctive screening support tool; its outputs are risk
indicators rather than medical diagnoses.
Formal IRB-approved clinical studies on a verified patient cohort are identified
as the necessary next step.

\textbf{Heuristic signal sources.}
All three score vectors ($\mathbf{p}^Q$, $\mathbf{p}^T$, $\mathbf{p}^A$) are
constructed from manually specified design constants rather than learned parameters.
While the audio branch extracts real-time physical acoustic parameters
(RMS energy, energy variability, zero-crossing rate, spectral centroid, spectral
rolloff, and speaking ratio) via the browser Web Audio API, its downstream mapping
to depression severity relies on a heuristic composite index rather than a deep
acoustic model trained on clinical recordings.
Future work will explore lightweight on-device ONNX/WebAssembly neural feature
extractors.

\textbf{Linguistic scope boundaries.}
While deterministic negation, attribution, quotation, and idiom rules have been integrated,
detecting subtle metaphorical, sarcastic, or culturally implicit distress
remains a limitation of rule-based filters.
Future work will evaluate lightweight zero-shot NLI models for deeper semantic
contextualization.

\textbf{Fusion weight optimization.}
The fusion weights ($w_T=0.50$, $w_A=0.30$, $w_Q=0.20$) are design choices.
While a five-scheme sensitivity analysis across constructed profiles was
conducted in Section~\ref{sec:eval}, formal optimization and empirical ablation
of fusion weights against ground-truth labels on an annotated clinical cohort
remain essential future work.

% ─────────────────────────────────────────────────────────────────────────────
\balance
\section{Conclusion}
\label{sec:conclusion}

This paper presented MindScreen, an engineering prototype of a full-stack
multimodal mental-health screening support system built around three principles:
decision-level late fusion, a formally specified
deterministic HRE module that prevents high-risk indicators from being diluted
by averaging, and a reproducible deterministic evidence script.
The HRE module produced the expected rule-compliant output across four isolated
constructed scenarios. In-process latency was measured for local deterministic
functions; hosted API, HTTP, database, browser, concurrent-load, and end-to-end
deployment latency remain unmeasured.
The system is explicitly presented as an engineering prototype; no clinical
diagnostic efficacy is claimed, and formal clinical evaluation remains the
prerequisite for healthcare deployment.

% ─────────────────────────────────────────────────────────────────────────────
\begin{thebibliography}{00}

\bibitem{who2023}
World Health Organization, ``Depressive disorder (depression),''
\textit{WHO Fact Sheets}, Mar.\ 2023.

\bibitem{gururaj2016}
G.~Gururaj \textit{et al.},
``National Mental Health Survey of India, 2015--16,''
\textit{NIMHANS Publication}, no.~129, 2016.

\bibitem{kessler2003}
R.~C.~Kessler \textit{et al.},
``Screening for serious mental illness in the general population,''
\textit{Arch.\ Gen.\ Psychiatry}, vol.~60, no.~2, pp.~184--189, 2003.

\bibitem{kroenke2001}
K.~Kroenke, R.~L.~Spitzer, and J.~B.~W.~Williams,
``The PHQ-9: Validity of a brief depression severity measure,''
\textit{J.\ Gen.\ Intern.\ Med.}, vol.~16, no.~9, pp.~606--613, 2001.

\bibitem{ji2022}
S.~Ji \textit{et al.},
``MentalBERT: Publicly available pretrained language models for mental
healthcare,''
in \textit{Proc.\ LREC}, 2022, pp.~7184--7190.

\bibitem{xu2024}
X.~Xu \textit{et al.},
``Mental-LLM: Leveraging large language models for mental health prediction,''
in \textit{Proc.\ ACL}, 2024, pp.~5120--5135.

\bibitem{cummins2015}
N.~Cummins \textit{et al.},
``A review of depression and suicide risk assessment using speech analysis,''
\textit{Speech Commun.}, vol.~71, pp.~10--49, 2015.

\bibitem{scherer2013}
S.~Scherer \textit{et al.},
``Automatic behavior descriptors for psychological disorder analysis,''
in \textit{Proc.\ IEEE FG}, 2013, pp.~1--8.

\bibitem{yang2023}
L.~Yang, D.~Jiang, and E.~Cambria,
``A survey on multimodal depression detection,''
\textit{IEEE Trans.\ Affect.\ Comput.}, vol.~14, no.~4,
pp.~3125--3144, 2023.

\bibitem{ringeval2019}
F.~Ringeval \textit{et al.},
``AVEC 2019 workshop and challenge,''
in \textit{Proc.\ ACM MM Workshop}, 2019, pp.~3--12.

\bibitem{devault2014}
T.~DeVault \textit{et al.},
``SimSensei Kiosk: A virtual human interviewer for healthcare decision
support,''
in \textit{Proc.\ AAMAS}, 2014, pp.~1061--1068.

\bibitem{fitzpatrick2017}
K.~K.~Fitzpatrick, A.~Darcy, and M.~Vierhile,
``Delivering cognitive behavior therapy to young adults with symptoms of
depression and anxiety using a fully automated conversational agent (Woebot),''
\textit{JMIR Ment.\ Health}, vol.~4, no.~2, p.~e19, 2017.

\bibitem{hartmann2022}
J.~Hartmann,
``Emotion English DistilRoBERTa-base,''
HuggingFace Model Hub, 2022.
[Online]. Available:
\url{https://huggingface.co/j-hartmann/emotion-english-distilroberta-base}

\bibitem{guo2017}
C.~Guo, G.~Pleiss, Y.~Sun, and K.~Q.~Weinberger,
``On calibration of modern neural networks,''
in \textit{Proc.\ ICML}, 2017, pp.~1321--1330.

\bibitem{telemanas2022}
Ministry of Health and Family Welfare, Government of India,
``Tele-MANAS,''
\textit{National Health Mission}, 2022.
[Online]. Available: \url{https://telemanas.mohfw.gov.in}

\bibitem{kiran2020}
Ministry of Social Justice and Empowerment, Government of India,
``KIRAN: 24/7 Mental Health Rehabilitation Helpline,''
2020.
[Online]. Available: \url{http://disabilityaffairs.gov.in}

\end{thebibliography}

\end{document}
